\chapter{线性系统的能控性和能观性}
\label{sec:04-controllability-observability}

\section{线性连续系统的能控性}

给定状态方程 $\dot x=Ax+Bu$，能控性问是否能在有限时间内把任意初态送至任意终态。教材也常从控制到原点定义能控性；对连续线性系统，状态转移矩阵可逆，所以这一定义与从原点到达任意终态等价。对线性定常系统，
\[
    x(t_f)=e^{A(t_f-t_0)}x(t_0)+\int_{t_0}^{t_f}e^{A(t_f-\tau)}Bu(\tau)\,\mathrm{d}\tau.
\]
由 Cayley--Hamilton 定理，$e^{A(t_f-\tau)}$ 可写成 $I,A,\dots,A^{n-1}$ 的线性组合。所有可由输入改变的终态方向组成 $\operatorname{im}Q_c$，其中
\[
    Q_c=[B\ AB\ \cdots\ A^{n-1}B],\qquad
    \text{完全能控}\iff\rank Q_c=n.
\]
若只要求控制输出 $y=Cx+Du$，相应的方向由
$Q_y=[CB\ CAB\ \cdots\ CA^{n-1}B\ D]$ 张成。此时满秩要求为 $\rank Q_y=p$，其中 $p$ 为输出维数。

\subsection{代数判据与模态判据}

PBH 判据检查 $A$ 每个特征值 $\lambda$ 是否满足 $\rank[\lambda I-A\ \ B]=n$。若存在非零左特征向量 $w^\top A=\lambda w^\top$ 且 $w^\top B=0$，则 $z=w^\top x$ 满足 $\dot z=\lambda z$；这个模态无法通过输入改变。当 $A$ 是上超对角线为 $1$ 的约旦形时，对同一特征值取各约旦块对应的 $B$ 的末行，它们必须线性无关。比如二阶约旦链的输入需能作用到链尾，才可沿超对角线向链首传递。

\ex{讲义思考\quad 状态变量与系统结构}{
引入状态变量后，需要回答哪些问题？为什么用经典控制理论的传递函数描述时，这些问题通常不显式出现？
}
\nt{\textbf{知识点}\quad 内部状态与外部输入输出描述，能控性、能观性和零极点相消。}
\sol{

状态方程把系统内部随时间演化的变量写了出来。要用输入控制这些变量，需要判断哪些状态方向能被输入作用到。要从传感器输出估计状态，需要判断哪些初态能被输入输出记录区分。能控性和能观性分别回答这两个问题，并决定能否任意配置全部状态反馈极点、能否根据输出构造状态观测器。

经典的传递函数从零初态出发描述输入到输出的关系。不能控模态无法由输入激发，不能观模态的变化无法传到输出，这些模态可以从约分后的传递函数中消失。比如
\[
    A=\begin{bmatrix}-1&0\\0&1\end{bmatrix},\qquad
    B=\begin{bmatrix}1\\0\end{bmatrix},\qquad C=[1\ 0],\qquad D=0
\]
给出 $G(s)=1/(s+1)$，而 $x_2(t)=e^t x_2(0)$。因此只观察 $G(s)$ 看不到第二个状态方向及其内部不稳定性。
}

\ex{讲义例题\quad 电桥的状态能控性（第9--11页）}{
下图中，电源电压 $u(t)$ 为输入，两个电容的端电压 $x_1(t),x_2(t)$ 为状态，正极性见图。设所有电阻和电容的参数为正。分析电桥平衡与不平衡时，输入对两个状态的控制能力。

\begin{center}
\begin{circuitikz}
    \draw (0,0) to[open,o-o] (0,3)
        to[C,l_=$C_1$] (2.55,3) -- (5,3)
        to[R,l^=$R_2$] (5,1.5) to[R,l^=$R_4$] (5,0) -- (0,0);
    \draw (2.55,3) to[R,l_=$R_1$] (2.55,1.5)
        to[R,l_=$R_3$] (2.55,0);
    \draw (2.55,1.5) to[C,l_=$C_2$] (5,1.5);
    \foreach \P in {(2.55,3),(2.55,1.5),(5,1.5),(2.55,0)} {
        \fill \P circle[radius=1.2pt];
    }
    \node[above] at (1.275,3.27) {$+\quad x_1\quad -$};
    \node[above] at (3.775,1.77) {$+\quad x_2\quad -$};
    \draw[->] (-.35,.2) -- node[left] {$u(t)$} (-.35,2.8);
\end{circuitikz}
\end{center}
}

\nt{\textbf{知识点}\quad 电桥平衡、电容电压作为状态、不能控方向的自由衰减、二阶能控矩阵。}

\sol 以底边为参考地。设电桥顶端电压为 $v=u-x_1$，左右中点为 $v_a,v_b$，故 $x_2=v_a-v_b$。记 $g_i=1/R_i$、$h=g_1+g_3$、$j=g_2+g_4$、$S=h+j$。两中点的电流方程相加，$C_2$ 的电流抵消，得到
\[
    h v_a+j v_b=(g_1+g_2)v .
\]
与 $v_a-v_b=x_2$ 联立后，可求两个电容的电流。
\[
\begin{aligned}
    C_1\dot x_1&=g_1(v-v_a)+g_2(v-v_b)=\alpha v-\beta x_2,\\
    C_2\dot x_2&=g_1v-hv_a=\beta v-\gamma x_2,\\
    \alpha&=\frac{(g_1+g_2)(g_3+g_4)}{S},\qquad
    \beta=\frac{g_1g_4-g_2g_3}{S},\qquad \gamma=\frac{hj}{S}.
\end{aligned}
\]
代入 $v=u-x_1$ 可读出 $A,B$。平衡条件为 $g_1g_4=g_2g_3$，即 $\beta=0$。此时 $\dot x_2=-(\gamma/C_2)x_2$，改变 $u$ 无法改变这一分量的运动。特别地，当 $R_1=R_2=R_3=R_4=R$ 时，
\[
    \dot x_1=\frac{u-x_1}{RC_1},\qquad
    \dot x_2=-\frac{x_2}{RC_2}.
\]
如果 $x_2(0)\ne0$，指数衰减只能使它逐渐趋零，任何有限时刻都无法靠 $u$ 把它带到零。

不平衡时 $\beta\ne0$。将 $B=[\alpha/C_1\ \ \beta/C_2]^\top$ 与 $AB$ 并列可得
\[
    \det[B\ AB]=-\frac{\beta(\alpha\gamma-\beta^2)}{C_1 C_2^2},\qquad
    \alpha\gamma-\beta^2=\frac{g_1g_2g_3+g_1g_2g_4+g_1g_3g_4+g_2g_3g_4}{S}>0.
\]
因此不平衡时秩为 $2$，系统完全能控。

\newcommand{\autonomousModeDiagram}[1]{%
\begin{tikzpicture}[>=Stealth,
    block/.style={draw,fill=mytheorembg,minimum width=10mm,minimum height=6mm},
    sum/.style={draw,circle,minimum size=5mm,inner sep=0pt}]
    \node[block] (gain1) at (0,0) {$-1$};
    \node[block] (int1) at (1.5,0) {$1/s$};
    \draw[->] (gain1) -- (int1);
    \draw[->] (int1) -- (3,0) node[right] {#1};
    \fill (2.28,0) circle[radius=1.15pt] node[above] {$x_1$};
    \draw[->] (2.28,0) -- ++(0,.65) -| (gain1.north);
    \node[sum] (sum2) at (3.35,-1.02) {$\scriptscriptstyle\sum$};
    \node[block] (int2) at (4.9,-1.02) {$1/s$};
    \node[block] (gain2) at (4.65,-1.86) {$-2$};
    \draw[->] (2.28,0) |- (sum2);
    \draw[->] (4.06,-.30) node[right] {$u$} -| (sum2.north);
    \draw[->] (sum2) -- (int2);
    \draw[->] (int2) -- (6.5,-1.02) node[right] {$x_2$};
    \fill (5.74,-1.02) circle[radius=1.15pt];
    \draw[->] (5.74,-1.02) |- (gain2);
    \draw[->] (gain2) -| (sum2);
\end{tikzpicture}%
}

\ex{讲义例题\quad 积分器结构中的自由状态（第12页）}{
图示系统满足 $\dot x_1=-x_1$、$\dot x_2=x_1-2x_2+u$。判断输入 $u$ 能否在有限时间内把任意初态转移到原点。
\begin{center}
\autonomousModeDiagram{}
\end{center}
}

\nt{\textbf{知识点}\quad 自由模态、有限时间转移以及能控子空间。}

\sol 第一积分器所在分支只有 $x_1$ 的自反馈。无论如何选择 $u$，始终有 $x_1(t)=e^{-t}x_1(0)$。只要 $x_1(0)\ne0$，这一状态在有限时刻都不为零。按矩阵检验，$B=[0\ \ 1]^\top$、$AB=[0\ \ {-2}]^\top$，能控矩阵秩为 $1$。能由输入改变的是 $x_2$ 方向，整个系统不完全能控。

\ex{讲义例题\quad 两个分量都受输入作用时的能控性（第13页）}{
给定
\[
    \dot x_1=-2x_1+x_2+u,\qquad \dot x_2=x_1-2x_2+u.
\]
判断两个状态能否用同一个输入在有限时间内同时置零。
}

\nt{\textbf{知识点}\quad 和差坐标、共模输入、能控性与状态坐标的关系。}

\sol 令 $z_+=x_1+x_2$、$z_-=x_1-x_2$。原方程相加与相减分别给出
\[
    \dot z_+=-z_++2u,\qquad \dot z_-=-3z_- .
\]
$u$ 只作用于和坐标。如果 $x_1(0)\ne x_2(0)$，则任意有限 $T$ 时都有 $z_-(T)=e^{-3T}z_-(0)\ne0$，故两个分量无法同时置零。检验原坐标下的矩阵，$B=[1\ \ 1]^\top$、$AB=-B$，秩仍为 $1$。


\ex{讲义例 4-1\quad 三阶友矩阵的能控性}{
判断系统的状态能控性，其中 $a_1,a_2,a_3$ 为任意实数。
\[
   \dot x=\begin{bmatrix}0&1&0\\0&0&1\\-a_3&-a_2&-a_1\end{bmatrix}x+
          \begin{bmatrix}0\\0\\1\end{bmatrix}u .
\]
}
\nt{\textbf{知识点}\quad 按 $B,AB,A^2B$ 构造能控矩阵，利用反三角结构判秩。}
\sol{
由 $AB=[0\ 1\ {-a_1}]^\top$、$A^2B=[1\ {-a_1}\ a_1^2-a_2]^\top$ 得
\[
    Q_c=\begin{bmatrix}0&0&1\\0&1&-a_1\\1&-a_1&a_1^2-a_2\end{bmatrix}.
\]
交换首末列后得到主对角元全为 $1$ 的下三角阵，故 $\det Q_c=-1$，任意参数下均完全能控。
}

\ex{讲义例 4-2、例 4-6\quad 秩判据与 PBH 判据}{
用能控矩阵和 PBH 判据分别分析下列系统的状态能控性（讲义第 23--24、36--37 页）。
\[
   \dot x=\begin{bmatrix}1&3&2\\0&2&0\\0&1&3\end{bmatrix}x+
          \begin{bmatrix}2&1\\1&1\\-1&-1\end{bmatrix}u .
\]
}
\nt{\textbf{知识点}\quad 代数判据给出能控子空间维数，PBH 判据给出失去控制的特征值。}
\sol{
计算得到
\[
    Q_c=\begin{bmatrix}
       2&1&3&2&5&4\\1&1&2&2&4&4\\-1&-1&-2&-2&-4&-4
    \end{bmatrix}.
\]
后二行互为相反数，前两列线性无关，故 $\rank Q_c=2$。再用 PBH 判据定位。$\det(\lambda I-A)=(\lambda-1)(\lambda-2)(\lambda-3)$，三个左特征向量可分别取
\[
    w_1^\top=[1\ {-2}\ {-1}],\qquad
    w_2^\top=[0\ 1\ 0],\qquad w_3^\top=[0\ 1\ 1].
\]
它们左乘 $B$ 分别得到 $[1\ 0]$、$[1\ 1]$、$[0\ 0]$。因此仅 $\lambda=3$ 对应的模态不能控。原方程中也可直接求和得到 $\dot x_2+\dot x_3=3(x_2+x_3)$。
}

\ex{讲义例 4-3\quad 含重根约旦形的能控性}{
分别判断 $\dot x=A_ix+B_iu$ 的状态能控性。
\[
\begin{aligned}
 (1)\ & A_1=\operatorname{diag}(-7,-5),\quad B_1=\begin{bmatrix}2\\5\end{bmatrix},\\
 (2)\ & A_2=\begin{bmatrix}-4&1&0\\0&-4&0\\0&0&-3\end{bmatrix},\quad
        B_2=\begin{bmatrix}0&0\\0&0\\1&1\end{bmatrix},\\
 (3)\ & A_3=\begin{bmatrix}-4&1&0&0\\0&-4&0&0\\0&0&-3&0\\0&0&0&-4\end{bmatrix},\quad
        B_3=\begin{bmatrix}0&0\\0&1\\2&0\\2&1\end{bmatrix},\\
 (4)\ & A_4=A_3,\quad B_4=\begin{bmatrix}0\\1\\2\\3\end{bmatrix}.
\end{aligned}
\]
}
\nt{\textbf{知识点}\quad 每个特征值分别收集约旦链尾的输入行，检查行向量是否独立。}
\sol{
（1）两个特征值各有一个一阶块，$B_1$ 的两行分别非零，完全能控。

（2）$\lambda=-4$ 的二阶块对应链尾为状态 $x_2$，其输入行全零。$x_1,x_2$ 构成无输入自治块，整体不完全能控。

（3）$\lambda=-4$ 有两个块，对应的链尾分别是 $x_2,x_4$，取到 $[0\ 1]$ 与 $[2\ 1]$，行列式 $-2\ne0$。$\lambda=-3$ 的链尾行 $[2\ 0]$ 非零，故完全能控。

（4）$\lambda=-4$ 需要在一个输入通道上独立作用于两个链尾，两个一维行向量 $[1]$ 与 $[3]$ 线性相关。非零模态组合 $x_4-3x_2$ 的导数为 $-4(x_4-3x_2)$，故整体不完全能控。
}

\ex{习题课 1\quad 状态能控性}{
判断以下两个系统的状态能控性。
\[
\begin{aligned}
 (1)\quad &\dot x=\begin{bmatrix}1&1\\1&0\end{bmatrix}x
             +\begin{bmatrix}0\\1\end{bmatrix}u,\\[2pt]
 (2)\quad &\dot x=\begin{bmatrix}-3&1&0\\0&-3&0\\0&0&-1\end{bmatrix}x
             +\begin{bmatrix}1&-1\\0&0\\2&0\end{bmatrix}
             \begin{bmatrix}u_1\\u_2\end{bmatrix}.
\end{aligned}
\]
}
\nt{\textbf{知识点}\quad 能控矩阵的秩判据，不能控模态与状态方程的对应关系。}
\sol{

（1）$AB=[1\ 0]^\top$，所以
\[
    Q_c=\begin{bmatrix}0&1\\1&0\end{bmatrix},\qquad
    \det Q_c=-1 .
\]
$Q_c$ 满秩，系统\textbf{状态完全能控}。输入先进入 $\dot x_2$，再经由 $x_2$ 影响 $\dot x_1$。

（2）逐次左乘 $A$ 后得到
\[
    Q_c=\begin{bmatrix}
    1&-1&-3&3&9&-9\\
    0&0&0&0&0&0\\
    2&0&-2&0&2&0
    \end{bmatrix}.
\]
第二行全为零，且 $B$ 的前两列线性无关，故 $\rank Q_c=2<3$。系统\textbf{状态不完全能控}。具体而言，$\dot x_2=-3x_2$，$x_2(t)$ 由初值唯一决定。
}

\ex{习题课 2\quad 输出能控性}{
判断下列系统的输出能控性。
\[
    \dot x=\begin{bmatrix}0&1&0\\0&0&1\\-6&-11&-6\end{bmatrix}x
          +\begin{bmatrix}0\\0\\1\end{bmatrix}u,\qquad
    y=\begin{bmatrix}1&0&0\end{bmatrix}x .
\]
}
\nt{\textbf{知识点}\quad 输出能控矩阵 $[CB\ CAB\ CA^2B\ D]$ 的秩与输出维数比较。}
\sol{

依次计算 $AB=[0\ 1\ {-6}]^\top$、$A^2B=[1\ {-6}\ 25]^\top$。$C$ 取第一行分量，且 $D=0$，故
\[
    Q_y=\begin{bmatrix}0&0&1&0\end{bmatrix},\qquad \rank Q_y=1 .
\]
该秩等于输出维数 $1$，系统\textbf{输出完全能控}。由 $\dot x_1=x_2$、$\dot x_2=x_3$ 也可看出，输入沿着 $x_3\to x_2\to x_1$ 传到输出。
}

\ex{讲义例 4-7\quad 输出能控与状态能控}{
判断 $\dot x=[1\ 1]^\top u,\ y=[1\ 1]x$ 的输出能控性（讲义第 42 页），并比较其状态能控性。
}
\nt{\textbf{知识点}\quad $Q_y$ 的秩与输出维数比较，$Q_c$ 的秩与状态维数比较。}
\sol{
这里 $A=0$、$D=0$，$CB=2$，故 $Q_y=[2\ 0\ 0]$ 的秩为 $1$，等于输出维数。取常值 $u=(y_f-y_0)/(2T)$，可在 $T>0$ 内把输出送至 $y_f$。与此同时，$Q_c=[B\ 0]$ 的秩为 $1<2$，且 $\dot x_1-\dot x_2=0$，因此状态不完全能控。
}

\subsection{时变系统的能控性}

设 $\Phi(t,\tau)$ 为由 $\tau$ 到 $t$ 的状态转移矩阵。定义终时刻坐标下的格拉姆矩阵
\[
    W_c=\int_{t_0}^{t_f}\Phi(t_f,\tau)B(\tau)B(\tau)^\top
                       \Phi(t_f,\tau)^\top\,\mathrm{d}\tau .
\]
该区间完全能控当且仅当 $W_c$ 非奇异。令 $d=x_f-\Phi(t_f,t_0)x(t_0)$，
输入 $u(\tau)=B(\tau)^\top\Phi(t_f,\tau)^\top W_c^{-1}d$ 代入状态解即得到期望终态。讲义也在初始时刻坐标下写格拉姆矩阵，两者通过可逆状态转移矩阵合同变换，有相同的秩。
若系数矩阵充分光滑，取 $M_0=B$，递推 $M_{k+1}=AM_k-\dot M_k$。在区间内某时刻，若 $[M_0\ \cdots\ M_{n-1}]$ 满行秩，则可判定能控。

\ex{讲义例 4-6\quad 时变系统在 $t_0=0$ 的能控性}{
判断下列时变系统在初始时刻 $t_0=0$ 的状态能控性（讲义第 47--50 页）。
\[
    \dot x(t)=\begin{bmatrix}0&t\\0&0\end{bmatrix}x(t)
                         +\begin{bmatrix}0\\1\end{bmatrix}u(t),\qquad t>0 .
\]
}
\nt{\textbf{知识点}\quad 时变能控格拉姆矩阵、导数秩充分条件。}
\sol{
齐次方程中 $x_2$ 为常数，将 $\dot x_1=t x_2$ 积分得
$\Phi(t,\tau)=\begin{bmatrix}1&(t^2-\tau^2)/2\\0&1\end{bmatrix}$。
按讲义在起点坐标计算，$\Phi(0,\tau)B=[-\tau^2/2\ 1]^\top$，故
\[
    W_c(0,T)=\begin{bmatrix}T^5/20&-T^3/6\\-T^3/6&T\end{bmatrix},\qquad
    \det W_c(0,T)=T^6/45>0\quad(T>0).
\]
所以每个长度为正的区间 $[0,T]$ 上完全能控。另一种验算是 $M_0=[0\ 1]^\top$、$M_1=[t\ 0]^\top$。$[M_0\ M_1]$ 在任何 $0<t<T$ 满秩，亦可判定能控。只在起始时刻 $t=0$ 检查会遗漏之后出现的耦合。
}

\section{线性连续系统的能观性}

给定 $u(t)$ 与 $y(t)$，能观性问能否在有限观测区间内唯一确定初态。先从输出中扣除已知输入引起的响应，所得 $\bar y(t)=Ce^{A(t-t_0)}x(t_0)$。在 $t_0$ 求 $0$ 至 $n-1$ 阶导数，得到
\[
    Q_o=\begin{bmatrix}C\\CA\\\vdots\\CA^{n-1}\end{bmatrix},\qquad
    \text{完全能观}\iff\rank Q_o=n.
\]
若两个初态之差属于 $\ker Q_o$，则在相同输入下它们的输出在所有时刻均相同。PBH 判据检查
$\rank\begin{bmatrix}\lambda I-A\\C\end{bmatrix}=n$；若 $Av=\lambda v$、$Cv=0$，则 $v$ 对应的模式无法观测。约旦形中，对同一特征值取各块对应的 $C$ 的首列，要检查这些列是否线性无关。

\ex{讲义例题\quad 根据电阻支路电流恢复电感电流（第53--56页）}{
如图，电压源 $u(t)$ 为输入，$R_3$ 支路的电流 $y=i_3$ 为输出。取 $x_1=i_1$、$x_2=i_2$，其中 $i_1$ 和 $i_2$ 分别是 $L_1$ 和 $L_2$ 的电流。分析在 $R_1=R_2$、$L_1=L_2$ 时能否由已知的输入和输出唯一确定状态，并讨论参数改变后如何判定。
\begin{center}
\begin{circuitikz}
    \draw (0,0) to[V,l_=$u(t)$] (0,2.75)
        to[R,l^=$R_1$] (1.88,2.75) to[L,l^=$L_1$,i>_=$i_1$] (4.3,2.75)
        -- (6.0,2.75) to[R,l^=$R_3$,i>_=$i_3$] (6.0,0) -- (0,0);
    \draw (4.3,2.75) to[R,l_=$R_2$] (4.3,1.4)
        to[L,l_=$L_2$,i>^=$i_2$] (4.3,0);
    \fill (4.3,2.75) circle[radius=1.2pt] (4.3,0) circle[radius=1.2pt];
    \node[left] at (-.05,1.93) {$+$};
    \node[left] at (-.05,.83) {$-$};
\end{circuitikz}
\end{center}
}

\nt{\textbf{知识点}\quad 节点电流、输出对称性、初态不可辨识、能观矩阵的秩。}

\sol 上方分流节点的电压记为 $v$。由节点电流定律，$y=x_1-x_2$，故 $v=R_3(x_1-x_2)$。两条含电感的回路给出
\[
    L_1\dot x_1=u-R_1x_1-v,\qquad
    L_2\dot x_2=v-R_2x_2 .
\]
当 $R_1=R_2=R,\ L_1=L_2=L$ 时，取 $u=0$、$x_1(t_0)=x_2(t_0)=a$。两支路的状态均为 $a e^{-(R/L)(t-t_0)}$，而 $y(t)$ 恒为零。任取不同的 $a$ 会得到同一条零输出，故状态不完全能观。

一般地，令 $r_i=R_i/L_i$、$q=R_3(1/L_1+1/L_2)$。由以上方程算得
\[
  C=\begin{bmatrix}1&-1\end{bmatrix},\qquad
  CA=\begin{bmatrix}-r_1-q&r_2+q\end{bmatrix},\qquad
  \det\begin{bmatrix}C\\CA\end{bmatrix}=r_2-r_1 .
\]
因此该电路完全能观的条件是 $R_1/L_1\ne R_2/L_2$。只改变一个元件的参数时应重新比较两个比值。

\ex{讲义例题\quad 开关断开后的电容电压能否观测（第57页）}{
下图中，$u(t)$ 是输入电压，右侧端口电压 $y(t)$ 是输出。取 $x_1=u_C$（左正右负）、$x_2=i_L$（向下为正）。开关 $K$ 在 $t_0$ 断开。判断断开后能否仅凭 $u,y$ 恢复两个状态。
\begin{center}
\begin{circuitikz}
    \draw (0,0) to[V,l_=$u(t)$] (0,1.4)
        to[nos,l_=$K$] (0,3) -- (1,3)
        to[R,l^=$R_1$] (3.35,3) -- (5.15,3);
    \draw (1,3) -- (1,2.03) to[C,l_=$C$] (3.35,2.03) -- (3.35,3);
    \node[above] at (2.175,2.24) {$+\quad u_C\quad -$};
    \draw (5.15,3) to[R,l^=$R_2$] (5.15,1.5)
        to[L,l^=$L$,i>_=$i_L$] (5.15,0) -- (0,0);
    \draw (5.15,3) -- (6.65,3) node[ocirc]{};
    \draw (5.15,0) -- (6.65,0) node[ocirc]{};
    \draw[->] (6.65,.25) -- node[right] {$y(t)$} (6.65,2.75);
    \node[left] at (-.05,1.29) {$+$};
    \node[left] at (-.05,.18) {$-$};
    \foreach \P in {(1,3),(3.35,3),(5.15,3),(5.15,0)} {
        \fill \P circle[radius=1.2pt];
    }
\end{circuitikz}
\end{center}
}

\nt{\textbf{知识点}\quad 开关拓扑变化、RC 自由放电、用两组不可区分初态证明能观性不足。}

\sol 断开开关后，$R_1$ 与 $C$ 的左侧公共节点没有外部电流。对该节点列电流方程得
\[
     C\dot x_1+x_1/R_1=0,\qquad
     x_1(t)=x_1(t_0^+)e^{-(t-t_0)/(R_1C)} .
\]
因此两条并联支路给右侧节点的净电流为零。取两组均与断开后理想电路相容的初态，令 $i_L(t_0^+)=0$，但令 $u_C(t_0^+)$ 不同。右侧支路电流及输出电压均可一直为零，电容仍按不同初值放电。由同一输入和相同输出无法反推 $u_C(t_0^+)$，故断开后的电容电压不可观。

\ex{讲义例题\quad 测得一个状态能否恢复另一个？（第58--59页）}{
已知
\[
     \dot x_1=-x_1,\qquad \dot x_2=x_1-2x_2+u,\qquad y=x_1.
\]
输入 $u$ 已知。根据图示模型，判断测量 $y$ 能否唯一确定 $x_1,x_2$。
\begin{center}
\autonomousModeDiagram{$y$}
\end{center}
}

\nt{\textbf{知识点}\quad 可直接测得的状态、未知初态的自由项、输出中缺失的状态方向。}

\sol 因为 $y=x_1$，第一状态随时可读。第二状态为
\[
     x_2(t)=e^{-2(t-t_0)}x_2(t_0)
        +\int_{t_0}^t e^{-2(t-\tau)}[y(\tau)+u(\tau)]\,d\tau .
\]
积分部分由已知信号决定。$x_2(t_0)$ 可以任取而不改变 $y$，所以第二状态无法由 $u,y$ 恢复，系统不完全能观。

\ex{讲义例题\quad 同一衰减速率下的和信号（第60--61页）}{
给定 $\dot x_1=-x_1$、$\dot x_2=-x_2$，只测量 $y=x_1+x_2$。问能否通过任意一段输出记录分别恢复 $x_1(t_0)$ 和 $x_2(t_0)$。
}

\nt{\textbf{知识点}\quad 能观性取决于输出随时间提供的独立信息。}

\sol 对两个标量状态分别积分，有 $y(t)=e^{-(t-t_0)}[x_1(t_0)+x_2(t_0)]$。因此测量记录只给出初态之和。初态 $[a\ \ b]^\top$ 与 $[a+\eta\ \ b-\eta]^\top$ 在每个时刻的输出完全相同。
$CA=-C$，能观矩阵秩为 $1$，不可观方向为 $[1\ \ {-1}]^\top$。


\ex{讲义例 4-7\quad 能观矩阵的秩}{
判断 $\dot x=\begin{bmatrix}-4&5\\1&0\end{bmatrix}x,\ y=[1\ {-1}]x$ 的状态能观性（讲义第 68 页）。
}
\nt{\textbf{知识点}\quad 如果 $CA$ 是 $C$ 的倍数，输出的导数不增加独立信息。}
\sol{
$CA=[-5\ 5]=-5C$，故 $\rank Q_o=1$，状态不完全能观。任取两组初态相差 $\alpha[1\ 1]^\top$，其中 $\alpha\ne0$，初始输出与之后的输出都相同，因为 $\dot y=-5y$。
}

\ex{讲义例 4-8\quad 约旦形的能观性}{
判断以下三个齐次系统 $\dot x=A_ix,\ y=C_ix$ 的状态能观性。
\[
\begin{aligned}
 (1)\ & A_1=\operatorname{diag}(-7,-5),\quad C_1=\begin{bmatrix}3&0\end{bmatrix},\\
 (2)\ & A_2=\begin{bmatrix}-4&1&0\\0&-4&0\\0&0&-3\end{bmatrix},\quad
        C_2=\begin{bmatrix}-1&2&1\end{bmatrix},\\
 (3)\ & A_3=\begin{bmatrix}-4&1&0&0\\0&-4&0&0\\0&0&-3&0\\0&0&0&-4\end{bmatrix},\quad
        C_3=\begin{bmatrix}1&1&0&-2\\0&2&1&0\end{bmatrix}.
\end{aligned}
\]
}
\nt{\textbf{知识点}\quad 每个特征值分别收集约旦链首的输出列。}
\sol{
（1）第二列为零，特征值 $-5$ 对应状态 $x_2$ 不可观。

（2）$\lambda=-4$ 的首列为 $-1$，$\lambda=-3$ 的首列为 $1$，两者均非零，系统完全能观。

（3）同一特征值 $\lambda=-4$ 对应的链首列为 $[1\ 0]^\top$ 与 $[-2\ 0]^\top$，线性相关。可取 $v=[2\ 0\ 0\ 1]^\top$ 检验 $A_3v=-4v$、$C_3v=0$。$\lambda=-3$ 的首列 $[0\ 1]^\top$ 非零。由于存在不可观初态方向，整体不完全能观。
}

\ex{讲义例 4-9\quad PBH 判据定位不可观模态}{
判断下列系统的状态能观性。
\[
    \dot x=\begin{bmatrix}0&1&0\\0&0&1\\-6&-11&-6\end{bmatrix}x,\qquad
    y=\begin{bmatrix}4&5&1\end{bmatrix}x .
\]
}
\nt{\textbf{知识点}\quad 通过 $Cv_\lambda$ 检查各右特征向量对应的模态。}
\sol{
特征多项式为 $(s+1)(s+2)(s+3)$，每个特征值 $\lambda$ 的右特征向量可取 $v_\lambda=[1\ \lambda\ \lambda^2]^\top$。计算 $Cv_\lambda=\lambda^2+5\lambda+4=(\lambda+1)(\lambda+4)$。$\lambda=-1$ 时为零，另外两个特征值代入均非零，故恰有 $-1$ 对应的模态不可观。若初态为 $[1\ {-1}\ 1]^\top$，输出恒为零。
}

\ex{习题课 3\quad 状态能观性}{
判断下列系统的状态能观性。
\[
    \dot x=\begin{bmatrix}0&1&0\\0&0&1\\-2&-4&-3\end{bmatrix}x,\qquad
    y=\begin{bmatrix}0&1&-1\\1&2&1\end{bmatrix}x .
\]
}
\nt{\textbf{知识点}\quad 多输出系统的能观矩阵，利用非零子式求秩。}
\sol{

将 $C,CA,CA^2$ 由上至下堆叠，得到
\[
    Q_o=\begin{bmatrix}
      0&1&-1\\1&2&1\\
      2&4&4\\-2&-3&-1\\
      -8&-14&-8\\2&2&0
    \end{bmatrix}.
\]
取前三行组成三阶方阵。第三行减去第二行的两倍后为 $[0\ 0\ 2]$，故该子式为 $2(0\cdot2-1\cdot1)=-2$。所以 $\rank Q_o=3$，系统\textbf{状态完全能观}。在这道题的齐次系统中，$y_1,y_2,\dot y_1$ 这三个标量关系就足以求出状态。
}

\subsection{时变系统的能观性}

若从输出中扣除已知输入的响应，所得 $\bar y(t)=C(t)\Phi(t,t_0)x(t_0)$。因此
\[
    W_o=\int_{t_0}^{t_f}\Phi(\tau,t_0)^\top C(\tau)^\top C(\tau)
                        \Phi(\tau,t_0)\,\mathrm{d}\tau .
\]
$W_o$ 非奇异当且仅当该区间完全能观。因为对于任一向量 $v$，$v^\top W_ov$ 等于假设初态差为 $v$ 时的输出差能量，$W_o$ 奇异即表示存在非零而不可辨认的初态差。若系数光滑，取 $N_0=C$，递推 $N_{k+1}=N_k A+\dot N_k$。区间内某时刻纵向堆叠 $N_0,\dots,N_{n-1}$ 得满列秩，也是完全能观的充分条件。

\ex{讲义例 4-10\quad 时变能观性}{
判断该系统在 $t_0=1$ 的状态能观性。
\[
    \dot x(t)=\begin{bmatrix}t&1&0\\0&t&0\\0&0&t^2\end{bmatrix}x(t),
    \qquad y(t)=\begin{bmatrix}1&0&1\end{bmatrix}x(t).
\]
}
\nt{\textbf{知识点}\quad 求输出的连续导数时，需同时对系数矩阵和状态求导。}
\sol{
$\dot y=(CA+\dot C)x$，依此递推得
\[
    \begin{bmatrix}N_0\\N_1\\N_2\end{bmatrix}
    =\begin{bmatrix}1&0&1\\t&1&t^2\\t^2+1&2t&t^4+2t\end{bmatrix}.
\]
行列式为 $t^2(t-1)^2+2t-1$，当 $t\ge1$ 时为正。观测区间内任一时刻的 $y,\dot y,\ddot y$ 可唯一确定该时刻的状态，状态转移矩阵可逆，故也确定 $x(1)$。系统在 $t_0=1$ 完全能观。
}

\section{对偶性原理}

$(A,B)$ 能控等价于 $(A^\top,B^\top)$ 能观，$(A,C)$ 能观等价于 $(A^\top,C^\top)$ 能控。其依据是
\[
    Q_o(A^\top,B^\top)=Q_c(A,B)^\top,\qquad
    Q_c(A^\top,C^\top)=Q_o(A,C)^\top .
\]
\begin{center}
\renewcommand{\arraystretch}{1.38}
\begin{tabular}{@{}lll@{}}
\toprule
& 状态能控 & 状态能观\\
\midrule
矩阵增广 & 横向拼接 $B, AB,\ldots$ & 纵向堆叠 $C, CA,\ldots$\\
代数判据 & $Q_c$ 满行秩 & $Q_o$ 满列秩\\
模态判据 & 每个左特征向量方向能受输入作用 & 每个右特征向量方向能反映在输出中\\
约旦形检查 & 同一特征值各块尾行线性无关 & 同一特征值各块首列线性无关\\
\bottomrule
\end{tabular}
\end{center}
对偶变换还可用来从能控规范形推导能观规范形，只需注意 $B,C$ 的角色和转置。

\section{结构分解与零极点相消}

设 $x=Pz$，则 $\tilde A=P^{-1}AP$、$\tilde B=P^{-1}B$、$\tilde C=CP$。$D$ 直接沿用。
\begin{itemize}
    \item 做\textbf{能控性分解}时，先从 $Q_c$ 中取一组极大线性无关列，放在 $P$ 的前几列，再补齐为可逆矩阵。$\operatorname{im}Q_c$ 对 $A$ 不变，因此 $\tilde A$ 左下块为零，$\tilde B$ 下块为零。取前一块得到能控子系统。
    \item 做\textbf{能观性分解}时，从 $Q_o$ 中取一组极大线性无关行，放在 $P^{-1}$ 的前几行，再补齐为可逆矩阵。$\ker Q_o$ 对 $A$ 不变，因此 $\tilde A$ 右上块为零，$\tilde C$ 右块为零。取前一块得到能观子系统。
    \item \textbf{能控能观分解}将状态分成能控能观、能控不能观、不能控能观、不能控不能观四部分。由零初态求传递函数时，只需保留能控能观部分。给定 $G(s)$ 的实现为最小实现，当且仅当该实现完全能控且完全能观。
\end{itemize}

\begin{note}
    $G(s)=C(sI-A)^{-1}B+D$。对 SISO 系统，将 $G$ 约分后，若分母次数小于 $\dim A$，说明实现中有不能控或不能观的模态。初态引起的自由响应仍应由原系统分析。
\end{note}

令 $\mathcal R=\operatorname{im}Q_c$、$\mathcal N=\ker Q_o$。$\mathcal R\cap\mathcal N$ 是能控但不能观的方向，最小实现的维数为
\[
    n_{\min}=\dim\mathcal R-\dim(\mathcal R\cap\mathcal N).
\]
图中分别画出两次分解允许的耦合。能控性分解后，不能控块可以通过 $A_{12}$ 影响能控块。提取能控子系统时令不能控块的初态为零。能观性分解后，不可观块受能观块影响，但它的初态不会影响输出。

\begin{center}
\begin{tikzpicture}[>=Latex, every node/.style={font=\small},
    block/.style={draw=whusakura,fill=whusakura!7,rounded corners=2pt,
      minimum width=2.5cm,minimum height=8mm,align=center}]
  \node[block] (c) at (0,0) {能控块 $z_c$};
  \node[block] (nc) at (0,1.35) {不能控块 $z_{\bar c}$};
  \draw[->,thick] (-2.3,0) node[left] {$u$} -- (c);
  \draw[->,thick] (nc) -- node[right] {$A_{12}$} (c);
  \draw[->,thick] (c.east) -- (2.3,0) node[right] {$y$};
  \draw[->,thick] (nc.east) -| (2.3,0);
  \node at (0,-.7) {能控性分解（示意）};
  \begin{scope}[xshift=6.4cm]
    \node[block] (o) at (0,0) {能观块 $z_o$};
    \node[block] (no) at (0,1.35) {不能观块 $z_{\bar o}$};
    \draw[->,thick] (-2.3,0) node[left] {$u$} -- (o);
    \draw[->,thick] (-2.3,0) |- (no);
    \draw[->,thick] (o) -- node[right] {$A_{21}$} (no);
    \draw[->,thick] (o.east) -- (2.3,0) node[right] {$y$};
    \node at (0,-.7) {能观性分解（示意）};
  \end{scope}
\end{tikzpicture}
\end{center}

\ex{讲义例 4-11\quad 能控性分解与自主演化块}{
求下列系统的能控子系统。
\[
    \dot x=\begin{bmatrix}1&2&-1\\0&1&0\\1&-4&3\end{bmatrix}x+
           \begin{bmatrix}0\\0\\1\end{bmatrix}u,\qquad y=[1\ {-1}\ 1]x.
\]
}
\nt{\textbf{知识点}\quad 先选能控列空间的基，分离自主块，区分完整变换与子系统。}
\sol{
$AB=[-1\ 0\ 3]^\top$，$A^2B=[-4\ 0\ 8]^\top=-4B+4AB$。按题面的 $A$ 计算，$A^2B$ 第三项为 $8$，与讲义第 103 页能控矩阵所列的 $0$ 有出入。前两列独立，故能控空间为二维。取 $P=[B\ AB\ e_2]$，则 $x=Pz$ 给出 $z_1=3x_1+x_3,\ z_2=-x_1,\ z_3=x_2$。因为 $Ae_2=[2\ 1\ {-4}]^\top$，它在新基底中的坐标为 $[2\ {-2}\ 1]^\top$。所以
\[
    \dot z=\begin{bmatrix}0&-4&2\\1&4&-2\\0&0&1\end{bmatrix}z+
          \begin{bmatrix}1\\0\\0\end{bmatrix}u,\qquad
    y=\begin{bmatrix}1&2&-1\end{bmatrix}z.
\]
$z_3$ 满足 $\dot z_3=z_3$。限制到 $z_3=0$，得
$\dot z_c=\begin{bmatrix}0&-4\\1&4\end{bmatrix}z_c+[1\ 0]^\top u$，
$y_c=[1\ 2]z_c$，其中 $z_c=[z_1\ z_2]^\top$。
}

\ex{习题课 4\quad 能控性结构分解}{
对下列系统进行能控性结构分解，写出能控子系统。
\[
    \dot x=\begin{bmatrix}1&1&0\\0&1&0\\0&1&1\end{bmatrix}x
          +\begin{bmatrix}0&1\\1&0\\0&1\end{bmatrix}u,\qquad
    y=\begin{bmatrix}1&1&1\end{bmatrix}x .
\]
}
\nt{\textbf{知识点}\quad 用能控子空间的基构造 $P$，计算新坐标下的矩阵并提取能控块。}
\sol{

设 $b_1=[0\ 1\ 0]^\top$、$b_2=[1\ 0\ 1]^\top$，由 $Ab_1=b_1+b_2$、$Ab_2=b_2$ 可知，$B$ 的两列张成的子空间对 $A$ 不变。因为 $b_1,b_2$ 线性无关，$\rank Q_c=2$。补入 $e_1=[1\ 0\ 0]^\top$，取
\[
    P=[b_1\ b_2\ e_1]=
    \begin{bmatrix}0&1&1\\1&0&0\\0&1&0\end{bmatrix},\qquad
    P^{-1}=\begin{bmatrix}0&1&0\\0&0&1\\1&0&-1\end{bmatrix}.
\]
令 $z=P^{-1}x$，得到 $z_1=x_2$、$z_2=x_3$、$z_3=x_1-x_3$。将坐标关系代回原方程得
\[
    \dot z=\begin{bmatrix}1&0&0\\1&1&0\\0&0&1\end{bmatrix}z
          +\begin{bmatrix}1&0\\0&1\\0&0\end{bmatrix}u,\qquad
    y=\begin{bmatrix}1&2&1\end{bmatrix}z .
\]
其中 $\dot z_3=z_3$。令 $z_c=[z_1\ z_2]^\top$，取前两维得到\textbf{能控子系统}
\[
    \dot z_c=\begin{bmatrix}1&0\\1&1\end{bmatrix}z_c
                    +\begin{bmatrix}1&0\\0&1\end{bmatrix}u,\qquad
    y_c=\begin{bmatrix}1&2\end{bmatrix}z_c .
\]
}

\ex{讲义例 4-12\quad 能观性分解（第二项输入系数为 $-1$）}{
求下列系统的能观子系统。
\[
    \dot x=\begin{bmatrix}0&0&-1\\1&0&-3\\0&1&-3\end{bmatrix}x+
           \begin{bmatrix}1\\-1\\0\end{bmatrix}u,\qquad
    y=[0\ 1\ {-2}]x .
\]
}
\nt{\textbf{知识点}\quad 能观坐标由 $C,CA$ 选取，能观子系统输入必须用本题的 $B$ 计算。}
\sol{
$CA=[1\ {-2}\ 3]$、$CA^2=-C-2CA$，故能观空间维数为 $2$。取 $z_1=Cx,\ z_2=CAx,\ z_3=x_3$，所得坐标变换可逆。按题面 $B=[1\ {-1}\ 0]^\top$ 计算，$CB=-1,\ CAB=3$。因此
\[
    \dot z_1=z_2-u,\qquad \dot z_2=-z_1-2z_2+3u,\qquad y=z_1 .
\]
不可观坐标满足 $\dot z_3=z_1-z_3$。讲义后续计算所列的 $[1\ {-1}]^\top$ 输入系数对应下一题的 $B=[1\ 1\ 0]^\top$。
}

\ex{习题课 5／讲义例 4-13\quad 能观与能控能观分解}{
给定系统
\[
    \dot x=\begin{bmatrix}0&0&-1\\1&0&-3\\0&1&-3\end{bmatrix}x
          +\begin{bmatrix}1\\1\\0\end{bmatrix}u,\qquad
    y=\begin{bmatrix}0&1&-2\end{bmatrix}x .
\]
\begin{enumerate}
    \item 进行能观性结构分解，写出能观子系统（习题课 5）。
    \item 按能控性与能观性进一步作联合结构分解（讲义例 4-13）。
\end{enumerate}
}
\nt{\textbf{知识点}\quad 用能观矩阵构造 $P^{-1}$，用能控矩阵构造 $P_c$，联合判断各分块的性质。}
\sol{

（1）计算 $CA=[1\ {-2}\ 3]$、$CA^2=[-2\ 3\ {-4}]=-C-2CA$。$C$ 与 $CA$ 线性无关，所以 $\rank Q_o=2$。用 $[0\ 0\ 1]$ 补齐，令
\[
    z=P^{-1}x,\qquad
    P^{-1}=\begin{bmatrix}0&1&-2\\1&-2&3\\0&0&1\end{bmatrix}.
\]
由第一、二行可得 $\dot z_1=CAx+CBu=z_2+u$，$\dot z_2=CA^2x+CABu=-z_1-2z_2-u$。由第三行可得 $\dot z_3=x_2-3x_3=z_1-z_3$，输出 $y=z_1$。因此 \textbf{能观子系统}为
\[
    \dot z_o=\begin{bmatrix}0&1\\-1&-2\end{bmatrix}z_o
             +\begin{bmatrix}1\\-1\end{bmatrix}u,\qquad
    y=\begin{bmatrix}1&0\end{bmatrix}z_o,\qquad
    z_o=\begin{bmatrix}z_1\\z_2\end{bmatrix}.
\]
这里 $z_3$ 可受 $z_1$ 驱动，而 $z_3$ 不出现在 $\dot z_o$ 和 $y$ 中。

（2）$B=[1\ 1\ 0]^\top$，$AB=[0\ 1\ 1]^\top$，$A^2B=-B-2AB$。取 $P_c=[B\ AB\ e_3]$，令 $v=P_c^{-1}x$。能控性分解得到
\[
    \dot v=\begin{bmatrix}0&-1&-1\\1&-2&-2\\0&0&-1\end{bmatrix}v+
          \begin{bmatrix}1\\0\\0\end{bmatrix}u,\qquad
    y=\begin{bmatrix}1&-1&-2\end{bmatrix}v .
\]
能控块的输出组合为 $v_1-v_2$，取 $w=[v_1-v_2\ \ v_2\ \ v_3]^\top$。计算得
\[
    \dot w=\begin{bmatrix}-1&0&1\\1&-1&-2\\0&0&-1\end{bmatrix}w+
          \begin{bmatrix}1\\0\\0\end{bmatrix}u,\qquad
    y=\begin{bmatrix}1&0&-2\end{bmatrix}w.
\]
$u$ 可达的子空间是 $w_3=0$。在此平面内，$w_1$ 进入输出，$w_2$ 无法影响输出；能观的两维子系统由 $w_1,w_3$ 构成，其能观阵为 $\begin{bmatrix}1&-2\\-1&3\end{bmatrix}$，行列式为 $1$。因此 $w_1$ 对应能控能观部分，$w_2$ 对应能控不能观部分，$w_3$ 对应不能控能观部分。不能控又不能观部分的维数为零。保留 $w_1$ 得最小实现 $\dot w_1=-w_1+u,\ y=w_1$，即 $G(s)=1/(s+1)$。
}

\ex{讲义例 4-14\quad 零极点相消}{
判断下列系统的传递函数是否有零极点相消，并求出约分结果。
\[
    \dot x=\begin{bmatrix}3&1\\2&2\end{bmatrix}x+
        \begin{bmatrix}1\\1\end{bmatrix}u,\qquad y=[1\ 0]x .
\]
}
\nt{\textbf{知识点}\quad 利用未约分的 $\det(sI-A)$ 定位被相消极点，用能控秩解释消去原因。}
\sol{
$AB=4B$，故 $Q_c$ 只有一维，$4$ 是唯一能控模态的特征值。$A$ 的特征多项式为 $(s-4)(s-1)$。直接用逆矩阵的第一行计算得到
\[
    G(s)=\frac{[s-2\ \ 1]\,[1\ 1]^\top}{(s-4)(s-1)}
        =\frac{s-1}{(s-4)(s-1)}=\frac{1}{s-4}.
\]
相消的极点为 $s=1$。全状态自由响应仍可能包含 $e^t$ 项。
}

\ex{纸质练习七\quad 串联系统的能控性、能观性和稳定性}{
设 $\boldsymbol{x}_1=[x_{11}\ x_{12}]^\top$，两个单输入单输出系统为
\[
\begin{aligned}
    S_1:\quad &\dot{\boldsymbol{x}}_1=
     \begin{bmatrix}0&1\\-3&-4\end{bmatrix}\boldsymbol{x}_1+
     \begin{bmatrix}0\\1\end{bmatrix}u_1,\qquad
     y_1=\begin{bmatrix}2&1\end{bmatrix}\boldsymbol{x}_1,\\
    S_2:\quad &\dot x_2=-2x_2+u_2,\qquad y_2=x_2 .
\end{aligned}
\]
按 $u_1\xrightarrow{\phantom{u}}\boxed{S_1}\xrightarrow{\,y_1=u_2\,}
\boxed{S_2}\xrightarrow{\phantom{y}}y_2$ 串联。
\begin{enumerate}
    \item 写出串联系统 $S$ 的状态方程和输出方程。
    \item 判定 $S$ 的能控性与能观性。若相应性质不满足，作结构分解并写出能控或能观子系统。
    \item 分析 $S$ 在 $x_e=0$ 的渐近稳定性与系统的 BIBO 稳定性。
\end{enumerate}
}
\nt{\textbf{知识点}\quad 串联建模，能控性分解，内部特征值与输入输出传递函数的稳定性判别。}
\sol{

（1）取 $u=u_1$，$y=y_2$，联合状态 $x=[x_{11}\ x_{12}\ x_2]^\top$。代入 $u_2=y_1=2x_{11}+x_{12}$，得
\[
    A=\begin{bmatrix}0&1&0\\-3&-4&0\\2&1&-2\end{bmatrix},\qquad
    B=\begin{bmatrix}0\\1\\0\end{bmatrix},\qquad
    C=\begin{bmatrix}0&0&1\end{bmatrix},\qquad D=0 .
\]

（2）并列 $B,AB,A^2B$，并堆叠 $C,CA,CA^2$，得到
\[
    Q_c=\begin{bmatrix}0&1&-4\\1&-4&13\\0&1&-4\end{bmatrix},\qquad
    Q_o=\begin{bmatrix}0&0&1\\2&1&-2\\-7&-4&4\end{bmatrix}.
\]
$Q_c$ 的首行、末行相同，且前两列无关，因此 $\rank Q_c=2$，$S$ 不能控。$\det Q_o=-1$，故 $S$ 完全能观。选取 $P=[B\ AB\ e_3]$，按 $z=P^{-1}x$ 换坐标。
\[
    P=\begin{bmatrix}0&1&0\\1&-4&0\\0&1&1\end{bmatrix},\qquad
    P^{-1}=\begin{bmatrix}4&1&0\\1&0&0\\-1&0&1\end{bmatrix}.
\]
变换后得到
\[
    \dot z=\begin{bmatrix}0&-3&0\\1&-4&0\\0&0&-2\end{bmatrix}z+
       \begin{bmatrix}1\\0\\0\end{bmatrix}u,\qquad
    y=\begin{bmatrix}0&1&1\end{bmatrix}z .
\]
取前两维即得能控子系统
\[
    \dot z_c=\begin{bmatrix}0&-3\\1&-4\end{bmatrix}z_c+
           \begin{bmatrix}1\\0\end{bmatrix}u,\qquad
    y_c=\begin{bmatrix}0&1\end{bmatrix}z_c .
\]
$z_3=x_2-x_{11}$ 满足 $\dot z_3=-2z_3$。该差值由初态决定，并在输出中出现，因此其模态不能控、能观。

（3）$A$ 为分块下三角矩阵，特征多项式为 $(s+1)(s+3)(s+2)$。三个特征值的实部均为负，所以原点\textbf{大范围渐近稳定}。从零初态计算
\[
    G_S(s)=G_2(s)G_1(s)
          =\frac{1}{s+2}\frac{s+2}{(s+1)(s+3)}
          =\frac{1}{(s+1)(s+3)}.
\]
输入输出通道中的 $-2$ 极点被相消。其冲激响应 $g(t)=\frac12(e^{-t}-e^{-3t})$ 在 $t\ge 0$ 上绝对可积，所以系统也\textbf{BIBO 稳定}。
}

\section{能控规范形和能观规范形}

设 $A$ 的特征多项式为 $s^n+a_{n-1}s^{n-1}+\cdots+a_0$。记 $F$ 为上超对角线为 $1$、末行为 $[-a_0\ \cdots\ -a_{n-1}]$ 的友矩阵，其余元素均为 $0$。记 $e_i$ 为第 $i$ 个单位列向量。本课所用四种形式为
\[
\begin{array}{c|c|c}
    & \text{I 形} & \text{II 形}\\ \hline
    (A_c,B_c) & (F^\top,e_1)&(F,e_n)\\
    (A_o,C_o) & (F,e_1^\top)&(F^\top,e_n^\top)
\end{array}
\]
对单输入能控系统，I 形可取 $P=Q_c$。求 II 形时，取 $\ell=e_n^\top Q_c^{-1}$，用 $\ell,\ell A,\ldots,\ell A^{n-1}$ 作为 $P^{-1}$ 的各行。对单输出能观系统，I 形可取 $P^{-1}=Q_o$。求 II 形时，取 $r=Q_o^{-1}e_n$，用 $r,Ar,\ldots,A^{n-1}r$ 作为 $P$ 的各列。

\ex{讲义例 4-15\quad 能控规范 I、II 形的变换}{
将 $\dot x=\begin{bmatrix}1&0\\-1&2\end{bmatrix}x+
\begin{bmatrix}-1\\1\end{bmatrix}u$ 分别化为能控规范 I 形和 II 形。
}
\nt{\textbf{知识点}\quad I 形用能控矩阵作基；II 形用 $Q_c^{-1}$ 的末行构造新坐标。}
\sol{
特征多项式为 $s^2-3s+2$，且 $Q_c=\begin{bmatrix}-1&-1\\1&3\end{bmatrix}$ 可逆。I 形取 $z=Q_c^{-1}x$。在基 $B,AB$ 中，$A B$ 的坐标为 $[0\ 1]^\top$，$A^2B=-2B+3AB$，所以
\[
    \dot z=\begin{bmatrix}0&-2\\1&3\end{bmatrix}z+
           \begin{bmatrix}1\\0\end{bmatrix}u .
\]
II 形取 $w_1=\ell x,\ w_2=\ell A x$，其中 $\ell=[1/2\ 1/2]$ 是 $Q_c^{-1}$ 的末行，且 $\ell A=[0\ 1]$。利用 $\ell B=0$、$\ell AB=1$，可得
\[
    \dot w=\begin{bmatrix}0&1\\-2&3\end{bmatrix}w+
           \begin{bmatrix}0\\1\end{bmatrix}u .
\]
}

\ex{讲义例 4-16\quad 能观规范 I、II 形的变换}{
将 $\dot x=\begin{bmatrix}1&-1\\0&2\end{bmatrix}x,\ y=[-1\ {-1/2}]x$
分别化为能观规范 I 形和 II 形。
}
\nt{\textbf{知识点}\quad I 形取 $z=Q_ox$；II 形用 $Q_o^{-1}$ 的末列构造变换矩阵。}
\sol{
$CA=[-1\ 0]$，$Q_o=\begin{bmatrix}-1&-1/2\\-1&0\end{bmatrix}$ 满秩，特征多项式为 $s^2-3s+2$。I 形中 $z_1=y$、$z_2=CAx$，$\dot z_1=z_2$，$\dot z_2=-2z_1+3z_2$，所以
\[
    \dot z=\begin{bmatrix}0&1\\-2&3\end{bmatrix}z,\qquad y=[1\ 0]z .
\]
II 形中，由 $Q_or=e_2$ 得 $r=[-1\ 2]^\top$，$Ar=[-3\ 4]^\top$。令 $x=[r\ Ar]w$。因为 $Cr=0$、$CAr=1$、$A^2r=-2r+3Ar$，于是
\[
    \dot w=\begin{bmatrix}0&-2\\1&3\end{bmatrix}w,\qquad y=[0\ 1]w .
\]
}

\ex{习题课 6\quad 能控规范 I 形和 II 形}{
将下列状态方程分别化为能控规范 I 形和 II 形。
\[
    \dot x=\begin{bmatrix}1&-2\\3&4\end{bmatrix}x+
                \begin{bmatrix}1\\1\end{bmatrix}u .
\]
}
\nt{\textbf{知识点}\quad 用 $Q_c$ 及其逆矩阵的末行构造规范形变换，特征多项式决定友矩阵中的系数。}
\sol{

先算特征多项式与能控矩阵。
\[
    \det(sI-A)=s^2-5s+10,\quad
    Q_c=\begin{bmatrix}1&-1\\1&7\end{bmatrix},\quad
    Q_c^{-1}=\frac{1}{8}\begin{bmatrix}7&1\\-1&1\end{bmatrix}.
\]
$Q_c$ 可逆，可以进行两种变换。

\textbf{I 形}取 $x=Q_c z$。因为 $Ab$ 是 $Q_c$ 的第二列，且由特征多项式有 $A^2b=-10b+5Ab$，
\[
    \dot z=\begin{bmatrix}0&-10\\1&5\end{bmatrix}z
               +\begin{bmatrix}1\\0\end{bmatrix}u .
\]

\textbf{II 形}取 $\ell=\frac18[-1\ 1]$、$\ell A=\frac18[2\ 6]$，令 $w_1=\ell x,\ w_2=\ell Ax$。由于 $\ell b=0$、$\ell Ab=1$，可写出 $\dot w_1=w_2$ 与 $\dot w_2=-10w_1+5w_2+u$，即
\[
    \dot w=\begin{bmatrix}0&1\\-10&5\end{bmatrix}w
               +\begin{bmatrix}0\\1\end{bmatrix}u .
\]
}

\ex{习题课 7\quad 能观规范 I 形和 II 形}{
将下列系统分别化为能观规范 I 形和 II 形。
\[
    \dot x=\begin{bmatrix}1&-1\\1&1\end{bmatrix}x+
                \begin{bmatrix}2\\1\end{bmatrix}u,\qquad
    y=\begin{bmatrix}-1&1\end{bmatrix}x .
\]
}
\nt{\textbf{知识点}\quad 能观规范形的坐标变换，变换时对 $B$ 和 $C$ 同时换坐标。}
\sol{

由 $CA=[0\ 2]$ 可得 $Q_o=\begin{bmatrix}-1&1\\0&2\end{bmatrix}$，行列式为 $-2$，能观。原矩阵的特征多项式为 $s^2-2s+2$。

\textbf{I 形}取 $z=Q_o x$。其中 $CB=-1$、$CAB=2$，再由 $A^2=-2I+2A$ 得
\[
    \dot z=\begin{bmatrix}0&1\\-2&2\end{bmatrix}z+
                \begin{bmatrix}-1\\2\end{bmatrix}u,\qquad
    y=\begin{bmatrix}1&0\end{bmatrix}z .
\]

\textbf{II 形}先用 $Q_o r=[0\ 1]^\top$ 解得 $r=[\tfrac12\ \tfrac12]^\top$，继而 $Ar=[0\ 1]^\top$。取 $x=Pw$，其中
\[
    P=[r\ Ar]=\begin{bmatrix}\tfrac12&0\\\tfrac12&1\end{bmatrix},\qquad
    P^{-1}=\begin{bmatrix}2&0\\-1&1\end{bmatrix}.
\]
应用 $P^{-1}AP,\ P^{-1}B,\ CP$，得到
\[
    \dot w=\begin{bmatrix}0&-2\\1&2\end{bmatrix}w+
                \begin{bmatrix}4\\-1\end{bmatrix}u,\qquad
    y=\begin{bmatrix}0&1\end{bmatrix}w .
\]
}

\section{传递函数矩阵的实现与最小实现}

由传递函数建立实现时，先通过多项式除法分离 $D$。若严格真部分为 $(b_0+b_1s+\cdots+b_{n-1}s^{n-1})/a(s)$，能控 II 形取 $C=[b_0\ b_1\ \cdots\ b_{n-1}]$。能观 II 形由其对偶得到，即将 $A$ 转置并交换 $B$ 与 $C$ 的转置。对传递矩阵，先将各通道的严格真部分化为同一分母，再用各阶分子系数矩阵构造分块规范形。

对真有理矩阵 $G(s)$，可用 $D=\lim_{s\to\infty}G(s)$ 求直接传递项。若严格真部分写为
\[
    \frac{B_0+B_1 s+\cdots+B_{n-1}s^{n-1}}{a(s)},\qquad
    a(s)=s^n+a_{n-1}s^{n-1}+\cdots+a_0 ,
\]
其中 $G$ 有 $p$ 个输出、$m$ 个输入，则 $B_i\in\mathbb R^{p\times m}$。构造能控 II 形时，按 $m$ 阶单位阵分块，状态维数为 $nm$。构造能观 II 形时，按 $p$ 阶单位阵分块，状态维数为 $np$。选择维数较小的构造后，还要验证另一项结构性质，或用独立的阶数下界证明最小性。

\ex{讲义例 4-17\quad 有直接传递项的 SISO 规范形实现}{
求
\[
    G(s)=\frac{2s^3+14s^2+26s+17}{s^3+6s^2+11s+6}
\]
的能控规范 I、II 形实现，并写出讲义用同一传递函数给出的能观规范 I、II 形实现。
}
\nt{\textbf{知识点}\quad 分出 $D$，用调换基底求 I 形，再由对偶性得到能观实现。}
\sol{
先做除法得到 $G(s)=2+(2s^2+4s+5)/a(s)$，其中 $a(s)=s^3+6s^2+11s+6$，因此四种形式的 $D=2$。记
$F=\begin{bmatrix}0&1&0\\0&0&1\\-6&-11&-6\end{bmatrix}$，$E=[5\ 4\ 2]$。能控 II 形为 $(F,e_3,E,2)$。取其能控矩阵 $P=[e_3\ Fe_3\ F^2e_3]$ 作变换，有
\[
    P=\begin{bmatrix}0&0&1\\0&1&-6\\1&-6&25\end{bmatrix},\qquad
    EP=[2\ {-8}\ 31].
\]
在该新基下，$P^{-1}FP=F^\top$，故能控 I 形为 $(F^\top,e_1,EP,2)$。由对偶性，能观 I 形为 $(F,(EP)^\top,e_1^\top,2)$，能观 II 形为 $(F^\top,E^\top,e_3^\top,2)$。
}

\ex{讲义例 4-18\quad 两输入一输出传递矩阵的最小实现}{
求
\[
    G(s)=\begin{bmatrix}\dfrac{1}{(s+1)(s+2)}&
                         \dfrac{1}{(s+2)(s+3)}\end{bmatrix}
\]
的最小实现。
}
\nt{\textbf{知识点}\quad 按输出数构造能观 II 形，利用跨通道出现的互异极点确定最小维数。}
\sol{
两个通道均严格真，$D=[0\ 0]$。公分母 $a(s)=(s+1)(s+2)(s+3)$，两个分子分别为 $s+3$ 和 $s+1$。取
\[
    A=\begin{bmatrix}0&0&-6\\1&0&-11\\0&1&-6\end{bmatrix},\quad
    B=\begin{bmatrix}3&1\\1&1\\0&0\end{bmatrix},\quad C=[0\ 0\ 1].
\]
由于 $C(sI-A)^{-1}=[1\ s\ s^2]/a(s)$，该实现复原了两个通道。$B$ 的前两列已张成前两个坐标的平面，而 $AB$ 的第一列含非零第三分量，故完全能控；能观 II 形本身完全能观。还可以由实际极点检验最小性：$-1,-2,-3$ 各自在至少一个通道出现，任何实现的特征多项式都须含这三个因子。

若分别实现两个二阶通道，同一极点 $-2$ 会被重复计入状态。公共分母的能观 II 形允许两路输入共享实现中的动态。
}

\ex{习题课 8\quad 传递函数的规范形实现}{
已知 $G(s)=\dfrac{s+4}{s^3+6s^2+11s+6}$，分别建立能控规范形和能观规范形实现。
}
\nt{\textbf{知识点}\quad 从分母构造友矩阵，按升幂排列严格真部分的分子系数，用对偶性得到能观形式。}
\sol{

分母按升幂的系数为 $6,11,6,1$，分子按升幂的系数为 $4,1,0$。$G$ 严格真，所以 $D=0$。能控 II 形为
\[
    A_c=\begin{bmatrix}0&1&0\\0&0&1\\-6&-11&-6\end{bmatrix},\quad
    B_c=\begin{bmatrix}0\\0\\1\end{bmatrix},\quad
    C_c=\begin{bmatrix}4&1&0\end{bmatrix}.
\]
令 $(sI-A_c)v=B_c$，前两行给出 $v_2=sv_1,\ v_3=s^2v_1$。第三行给出 $v_1=1/(s^3+6s^2+11s+6)$，于是 $C_cv=G(s)$。

将能控 II 形转置取对偶，得到能观 II 形
\[
    A_o=\begin{bmatrix}0&0&-6\\1&0&-11\\0&1&-6\end{bmatrix},\quad
    B_o=\begin{bmatrix}4\\1\\0\end{bmatrix},\quad
    C_o=\begin{bmatrix}0&0&1\end{bmatrix}.
\]
取对偶后传递函数转置，SISO 的传递函数仍为 $G(s)$。

若要求 I 形，可对上面的能控实现取 $P=[B_c\ A_cB_c\ A_c^2B_c]$，算得 $C_cP=[0\ 1\ {-2}]$。因此能控 I 形取 $(A_o,e_1,[0\ 1\ {-2}],0)$，其对偶 $(A_c,[0\ 1\ {-2}]^\top,e_1^\top,0)$ 即为能观 I 形。
}

\ex{习题课 9\quad 传递矩阵的最小实现}{
设
\[
    G(s)=\begin{bmatrix}
      \dfrac{2(s+3)}{(s+1)(s+2)} & \dfrac{4(s+4)}{s+5}
    \end{bmatrix}.
\]
求该系统的最小实现。
}
\nt{\textbf{知识点}\quad 提取直接传递项，单输出系统的能观规范 II 形实现，利用极点证明阶数最小。}
\sol{

第二个通道可写为 $4-4/(s+5)$。提取 $D=[0\ 4]$ 后通分，得
\[
\begin{aligned}
    a(s)&=(s+1)(s+2)(s+5)=s^3+8s^2+17s+10,\\
    G(s)-D&=\frac{[2s^2+16s+30\quad -4s^2-12s-8]}{a(s)} .
\end{aligned}
\]
按分子系数从常数项至二次项排 $B$ 的两列，得
\[
    A=\begin{bmatrix}0&0&-10\\1&0&-17\\0&1&-8\end{bmatrix},\quad
    B=\begin{bmatrix}30&-8\\16&-12\\2&-4\end{bmatrix},\quad
    C=\begin{bmatrix}0&0&1\end{bmatrix},\quad
    D=\begin{bmatrix}0&4\end{bmatrix}.
\]
其传递矩阵为题设 $G(s)$，因为 $C(sI-A)^{-1}=[1\ s\ s^2]/a(s)$。三个互异的极点 $-1,-2,-5$ 均在至少一个约分后的通道出现。任意实现中，每个通道的极点都是 $\det(sI-A)$ 的根，所以状态维数至少为 $3$。以上实现恰有三个状态，故为\textbf{最小实现}。
}

\ex{纸质练习四\quad 最小实现}{
求传递矩阵
\[
    G(s)=\begin{bmatrix}
    \dfrac{s+1}{(s+2)(s+3)}&\dfrac{5s+4}{s+5}
    \end{bmatrix}
\]
的最小实现。
}
\nt{\textbf{知识点}\quad 先做多项式除法，再通分并按升幂排列分子系数。}
\sol{

由 $\dfrac{5s+4}{s+5}=5-\dfrac{21}{s+5}$，先取 $D=[0\ 5]$。两通道严格真部分的公分母为 $a(s)=(s+2)(s+3)(s+5)=s^3+10s^2+31s+30$。第一通道的分子为 $(s+1)(s+5)=s^2+6s+5$，第二通道为 $-21(s+2)(s+3)=-21s^2-105s-126$。因此能观规范 II 形实现为
\[
    A=\begin{bmatrix}0&0&-30\\1&0&-31\\0&1&-10\end{bmatrix},\quad
    B=\begin{bmatrix}5&-126\\6&-105\\1&-21\end{bmatrix},\quad
    C=\begin{bmatrix}0&0&1\end{bmatrix},\quad
    D=\begin{bmatrix}0&5\end{bmatrix}.
\]
将 $B$ 左乘 $C(sI-A)^{-1}=[1\ s\ s^2]/a(s)$ 即得两通道的严格真部分。两个通道合计有三个互异的实际极点 $-2,-3,-5$，从而任何实现至少三阶。此实现为三阶，达到最小维数。
}

\section{MATLAB 验算}

用 \texttt{rank(ctrb(A,B))}、\texttt{rank(obsv(A,C))} 得到能控矩阵和能观矩阵的秩，分别与 \texttt{size(A,1)} 比较。用 \texttt{minreal(ss(A,B,C,D))} 求消去冗余动态后的实现。数值秩和约简结果受容差影响，含参数或接近秩亏的题目可用符号运算复核。

\Needspace{8\baselineskip}

\ex{讲义例题\quad 用 \texttt{ctrb} 判断多输入系统（第187--188页）}{
已知三阶系统 $\dot x=Ax+Bu$，其中
\[
    A=\begin{bmatrix}-2&2&-1\\0&-2&0\\1&-4&0\end{bmatrix},\qquad
    B=\begin{bmatrix}0&0\\0&1\\1&0\end{bmatrix}.
\]
判别系统的状态能控性。
}

\nt{\textbf{知识点}\quad 多输入能控矩阵按整块排列，列数可以大于状态维数。}

\sol 将 $AB$ 和 $A^2B$ 逐块接到 $B$ 之后，得
\[
    Q_c=\begin{bmatrix}
    0&0&-1&2&2&-4\\ 0&1&0&-2&0&4\\ 1&0&0&-4&-1&10
    \end{bmatrix}.
\]
前三列组成的方阵行列式为 $1$，故 $\rank Q_c=3$，系统完全能控。
\begin{verbatim}
A = [-2 2 -1; 0 -2 0; 1 -4 0];
B = [0 0; 0 1; 1 0];
Qc = ctrb(A,B); rank(Qc) == size(A,1)
\end{verbatim}

\ex{讲义例题\quad 用 \texttt{obsv} 判断重特征值系统（第189--190页）}{
已知 $\dot x=Ax$、$y=Cx$，其中
\[
A=\begin{bmatrix}4&1&0&0\\0&4&0&0\\0&0&4&1\\0&0&0&4\end{bmatrix},\qquad
C=\begin{bmatrix}1&1&2&1\\1&2&2&0\end{bmatrix}.
\]
判别系统的状态能观性。
}

\nt{\textbf{知识点}\quad 约旦块、幂零矩阵降幂、能观矩阵零空间。}

\sol 写 $A=4I+N$，其中 $N^2=0$。每一块 $CA^k$ 都是 $C$、$CN$ 的线性组合，因此只需对下面的矩阵求秩。
\[
    \begin{bmatrix}C\\CN\end{bmatrix}
    =\begin{bmatrix}1&1&2&1\\1&2&2&0\\0&1&0&2\\0&1&0&2\end{bmatrix}.
\]
前三行线性无关，第四行等于第三行，故秩为 $3<4$。不可观子空间由 $[-2\ \ 0\ \ 1\ \ 0]^\top$ 张成。
\begin{verbatim}
A = [4 1 0 0; 0 4 0 0; 0 0 4 1; 0 0 0 4];
C = [1 1 2 1; 1 2 2 0];
Qo = obsv(A,C); rank(Qo)
\end{verbatim}

\ex{讲义例题\quad 用 \texttt{ctrbf} 分解不能控系统（第191--193页）}{
给定 $\dot x=Ax+Bu$、$y=Cx$，其中
\[
    A=\begin{bmatrix}1&2&-1\\0&1&0\\1&-4&3\end{bmatrix},\quad
    B=\begin{bmatrix}0\\0\\1\end{bmatrix},\quad
    C=\begin{bmatrix}1&-1&1\end{bmatrix}.
\]
判断系统能控性。如果不完全能控，作能控性分解并写出能控子系统。
}

\nt{\textbf{知识点}\quad 能控空间由 $B,AB,\ldots$ 张成，\texttt{ctrbf} 的变换约定为 $z=Tx$。}

\sol 能控矩阵为
\[
    Q_c=\begin{bmatrix}0&-1&-4\\0&0&0\\1&3&8\end{bmatrix}.
\]
前两列独立且第二行全零，故秩为 $2$。$\dot x_2=x_2$ 是无法由输入改变的状态。取 $z=[x_2\ \ {-x_1}\ \ x_3]^\top$，将不能控分量排在最前，得
\[
    \bar A=\begin{bmatrix}1&0&0\\-2&1&1\\-4&-1&3\end{bmatrix},\quad
    \bar B=\begin{bmatrix}0\\0\\1\end{bmatrix},\quad
    \bar C=\begin{bmatrix}-1&-1&1\end{bmatrix}.
\]
因此能控子系统的状态为 $[z_2\ \ z_3]^\top$，对应 $A_c=\begin{bmatrix}1&1\\-1&3\end{bmatrix}$、$B_c=[0\ \ 1]^\top$、$C_c=[-1\ \ 1]$。求原系统非零初值响应时仍须保留 $z_1$ 及它对其余状态和输出的影响。
\begin{verbatim}
A = [1 2 -1; 0 1 0; 1 -4 3]; B = [0;0;1]; C = [1 -1 1];
[Ab, Bb, Cb, T, k] = ctrbf(A,B,C);
r = sum(k);
idx = size(A,1)-r+1 : size(A,1);
Ac = Ab(idx,idx); Bc = Bb(idx,:); Cc = Cb(:,idx);
\end{verbatim}

\ex{讲义例题\quad 用 \texttt{obsvf} 分解不能观系统（第194--196页）}{
给定 $\dot x=Ax+Bu$、$y=Cx$，其中
\[
    A=\begin{bmatrix}0&0&-1\\1&0&-3\\0&1&-3\end{bmatrix},\quad
    B=\begin{bmatrix}1\\1\\0\end{bmatrix},\quad
    C=\begin{bmatrix}0&1&-2\end{bmatrix}.
\]
判断系统能观性。如果不完全能观，作能观性分解并写出能观子系统。
}

\nt{\textbf{知识点}\quad 用 $C,CA$ 定义可由测量恢复的坐标，按 \texttt{obsvf} 的分块顺序提取子系统。}

\sol $CA=[1\ \ {-2}\ \ 3]$，$CA^2=[-2\ \ 3\ \ {-4}]=-2CA-C$，故 $Q_o$ 秩为 $2$。取 $z=[x_3\ \ CAx\ \ Cx]^\top$，这三个线性组合给出行列式为 $1$ 的变换。由 $\dot z_2=CA^2x+CABu$、$\dot z_3=CAx+CBu$ 等式得到
\[
    \bar A=\begin{bmatrix}-1&0&1\\0&-2&-1\\0&1&0\end{bmatrix},\quad
    \bar B=\begin{bmatrix}0\\-1\\1\end{bmatrix},\quad
    \bar C=\begin{bmatrix}0&0&1\end{bmatrix}.
\]
能观子系统由 $(z_2,z_3)$ 构成，$A_o=\begin{bmatrix}-2&-1\\1&0\end{bmatrix}$、$B_o=[-1\ \ 1]^\top$、$C_o=[0\ \ 1]$。$z_1$ 的初值不影响输出。手算选择整数组合，MATLAB 选择正交变换时矩阵元素可以不同，两种坐标下的维数和输入输出行为一致。
\begin{verbatim}
A = [0 0 -1; 1 0 -3; 0 1 -3]; B = [1;1;0]; C = [0 1 -2];
[Ab, Bb, Cb, T, k] = obsvf(A,B,C);
r = sum(k);
idx = size(A,1)-r+1 : size(A,1);
Ao = Ab(idx,idx); Bo = Bb(idx,:); Co = Cb(:,idx);
\end{verbatim}

